% !PRINTED 470
\weberpage{470}

%<<CONT
numerischen Factoren abgesehen, wie sich aus (8) leicht ergiebt,
der Reihe nach
\[
h,\ g,\ h^2,\ hg,\ (g^2,h^3),\ (g^2,h^2g),
\]
worin $(g^2,h^3)$, $(g^2,h^2g)$ lineare homogene Ausdr\"ucke mit numeri\-schen
Co\"efficienten bedeuten.

Die Rechnung kann ebenso ausgef\"uhrt werden, wie in dem
obigen speciellen Falle. Zur Vereinfachung kann man $x_3=0$
setzen und erh\"alt immer noch Gleichungen genug zur Be\-stimmung
aller Co\"efficienten. Man findet so die Co\"efficienten
der Reihe nach:
\[
7\,\frac{1-\sqrt{-7}}{2}\,h,
\]
\[
-7\,\frac{1-\sqrt{-7}}{2}\,g,
\]
\[
-7\,(4+\sqrt{-7})\,h^2,
\]
\[
14\,(2+\sqrt{-7})\,hg,
\]
\[
-7\,\frac{5+\sqrt{-7}}{2}\,g^2 - 7\cdot\frac{7+3\sqrt{-7}}{2}\,h^3,
\]
\[
-g^2 + \frac{167-7\sqrt{-7}}{2}\,gh^2.
\]
Die Functionen $g, h$ sind gebrochene Invarianten, die nur
von den Verh\"altnissen $x_1:x_2:x_3$ abh\"angen, und wir k\"onnen
jede andere Invariante von derselben Eigenschaft
\gesperrt{rational} durch $g, h$ ausdr\"ucken. Denn stellen wir eine
solche Invariante als Quotienten zweier Formen gleichen Grades
ohne gemeinsamen Theiler dar, so m\"ussen Z\"ahler und Nenner
ganze Invarianten gleichen Grades sein, weil n\"amlich zwei in
einfachster Form dargestellte gebrochene Functionen der
Variablen $x$ nur dann einander gleich sein k\"onnen, wenn Z\"ahler
und Nenner einzeln bis auf constante Factoren einander gleich
sind. Hier k\"onnen nun Z\"ahler und Nenner nicht von ungeradem
Grade sein, weil sie sonst den gemeinsamen Factor $K$ h\"atten
(\S.\,126). Also sind Z\"ahler und Nenner rational durch $f, C, \Delta$
darstellbar; und wenn ein im Z\"ahler oder im Nenner vorkommendes
Glied
\weq{13}{f^a\,\Delta^b\,C^c,}
und $m$ der Grad von Z\"ahler und Nenner ist, so ist
\weq{14}{4a + 6b + 14c = m.}
